% Generated draft: requires mathematical and source review.
\documentclass[11pt]{article}
\usepackage[margin=25mm]{geometry}
\usepackage{fontspec}
\setmainfont{Latin Modern Roman}
\usepackage{amsmath,amssymb,mathtools,unicode-math}
\setmathfont{Latin Modern Math}
\usepackage{enumitem,needspace,xurl,hyperref}
\setlength{\parindent}{0pt}
\setlength{\parskip}{5pt}
\setlength{\emergencystretch}{3em}
\newcommand{\N}{\mathbb{N}}
\newcommand{\Z}{\mathbb{Z}}
\newcommand{\Q}{\mathbb{Q}}
\newcommand{\R}{\mathbb{R}}
\newcommand{\C}{\mathbb{C}}
\newcommand{\F}{\mathbb{F}}
\newcommand{\A}{\mathbb{A}}
\newcommand{\PP}{\mathbb{P}}
\newcommand{\E}{\mathbb{E}}
\newcommand{\eps}{\varepsilon}
\newcommand{\dd}{\,\mathrm{d}}
\begin{document}
\section*{UnsolvedMath ID: 30001790}
\textbf{Source problem code:} catalogue entry
\subsection{Definitions and mathematical statement}
Let \(\Gamma\) be a countable group. A measure-preserving action \(\alpha\) on a measure space \((X,\Sigma,\mu)\) assigns to each \(g\in\Gamma\) an invertible measurable transformation \(\alpha_g\), with measurable inverse, such that \(\alpha_e=\operatorname{id}_X\), \(\alpha_{gh}=\alpha_g\circ\alpha_h\), and \(\mu(\alpha_g^{-1}(E))=\mu(E)\) for every \(E\in\Sigma\). The supplied record does not specify the measure spaces allowed, whether transformations are identified modulo null sets, or whether the actions must satisfy additional hypotheses.

For each admissible action, write \(R_\alpha\) for the crossed-product ring intended by the source and \(S_\alpha\) for its regular closure in the intended ambient ring. A ring \(S\) is von Neumann regular if
\[
\forall s\in S\;\exists t\in S:\quad s=sts.
\]
This definition alone does not specify the regular-closure construction: the coefficient ring, ambient ring, and any involution or additional closure requirements must also be given. If \(\rho_\alpha\) is the intended rank function on \(S_\alpha\), inducing the rank pseudometric
\[
d_\alpha(s,t)=\rho_\alpha(s-t),
\]
then \(A_\alpha\) denotes the corresponding completion, after identifying elements at distance zero when necessary. The record does not specify \(\rho_\alpha\) or its normalization.

\textbf{Question, conditional on specifying these constructions.} Fix a nonempty class \(\mathcal C_\Gamma\) of admissible measure-preserving actions and a specified notion \(\mathcal E\) of isomorphism for their completed algebras. Which invariants \(I\) of actions are determined by the isomorphism class of \(A_\alpha\)? Explicitly, for which such invariants does
\[
\forall\alpha,\beta\in\mathcal C_\Gamma,\qquad
A_\alpha\cong_{\mathcal E}A_\beta
\ \Longrightarrow\ I(\alpha)=I(\beta)
\]
hold? Further, characterize the countable nonamenable groups \(\Gamma\) for which the completion is independent of the action, in the sense that
\[
\exists B\;\forall\alpha\in\mathcal C_\Gamma,\qquad
A_\alpha\cong_{\mathcal E}B.
\]
In particular, does any countable nonamenable group have this property?

The choices of \(\mathcal C_\Gamma\), \(R_\alpha\), \(S_\alpha\), \(\rho_\alpha\), and \(\mathcal E\) remain unresolved in the supplied record. Thus these formulas describe the quantifier structure of the supplied question, rather than a uniquely specified mathematical conjecture. In particular, ``the same algebra'' is not specified to mean an abstract algebra isomorphism, an involution-preserving isomorphism, or an isomorphism preserving rank or other distinguished structure.
\subsection{Short English statement}
The question asks what information about a measure-preserving action of a countable group can be recovered from the rank completion of the regular closure of its crossed-product ring, and which countable nonamenable groups could give isomorphic completions for every admissible action. The admissible actions, ring construction, rank, and meaning of isomorphism require clarification before this becomes a fully specified problem.
\subsection{Sources}
\begin{itemize}
\item Supplied local UnsolvedMath registry record: permanent ID \(30001790\), title \emph{Rank Completions of Crossed-Product Rings}, source problem code \texttt{catalogue entry}, in the collection \emph{Oberwolfach Reports Open Problems}. The record gives proposed year 2011 but names no individual proposer.
\item Supplied report citation: \emph{Finite-dimensional Approximations of Discrete Groups} (2011), Oberwolfach Reports, pp.~1429--1467, DOI \texttt{10.4171/owr/2011/26}. Supplied source URL and literature-source link labeled ``original question'': \url{https://doi.org/10.4171/owr/2011/26}. The background describes Oberwolfach Reports as documenting mathematical workshops and their open-problem discussions. No problem-specific page or named contributor is supplied.
\item Literature assessment dated exactly \texttt{2026-08-21}: the supplied curation labels the problem open and reports that no general identification for the stated class of groups was found. This is supplied dated literature triage derived from OWR curation, not an independently verified assessment; the record expressly calls for independent verification of time-sensitive status and cited claims.
\item The registry labels the statement \texttt{corrected\_verified} and reports a manual semantic, typographic, and LaTeX review asserting that the cleanup is self-contained and adds no unsupported source-dependent content. These are supplied editorial assertions; this draft does not independently verify them or the underlying report.
\end{itemize}
\end{document}
